\chapter{Design and Architecture}
\label{chap:design}

\section{Design Goals}

TensorStore was designed around four primary requirements. First, the framework must support multiple I/O backends and checkpoint formats under a consistent interface, so that backend decisions can be studied without rewriting the rest of the loading stack. Second, the framework must keep the loader logic separate from the measurement harness, so that instrumentation does not distort the measured path. Third, the design must make its concurrency model explicit, because queue depth, worker fanout, and batching policy are central to the performance story. Fourth, the framework must remain accessible from Python and from a real serving engine, so that backend-level findings can be tested beyond Rust-only microbenchmarks.

These requirements led to a layered architecture that separates format parsing, I/O execution, policy selection, and evaluation.

\section{System Architecture}

TensorStore is structured into three layers: a format layer, a backend layer, and an evaluation layer.

The format layer parses checkpoint metadata and produces a normalised loading workload. The key abstraction is a request-local view of the checkpoint: the format layer determines which file ranges must be read, but it does not commit to how those reads are scheduled. This separation allows the same logical workload to be executed by different backends while preserving the semantics visible to the caller.

The backend layer implements I/O execution. Each backend receives the same logical workload and executes it according to its own submission model. In the final design, backends are allowed to perform backend-local planning such as chunking large files or coalescing nearby ranges, but they must preserve request-local output semantics. This was an important design refinement: the project began with a stricter ``fixed request trace'' framing, but the final system recognises that practical loaders need controlled internal planning if they are to be evaluated fairly as real implementations rather than as artificially constrained sketches.

The evaluation layer contains the Rust profiling harness, the Python benchmarks, and the vLLM integration benchmarks. These are kept outside the core loading path so that the production implementation and the measurement logic remain distinct.

\section{Common Backend Interface}

All eager-loading backends implement the same public surface:

\begin{itemize}
\item \texttt{load(path)} for a full file,
\item \texttt{load\_batch(paths)} for multiple whole-file loads,
\item \texttt{load\_range(path, offset, len)} for one byte range, and
\item \texttt{load\_range\_batch(requests)} for many ranges.
\end{itemize}

This interface is deliberately narrow. It is expressive enough for both \texttt{SafeTensors} and \texttt{ServerlessLLM}, but small enough that backend differences remain interpretable.

\section{Adaptive Loading Policy}

The central design decision of the final system is that backend selection is handled as an adaptive heuristic rather than as a fixed policy. The project began as a comparison between synchronous I/O and \texttt{io\_uring}, but the final evidence showed that no single backend is best across all checkpoint sizes and layouts.

The final \texttt{default} policy therefore uses workload-aware signals. For \texttt{SafeTensors}, the key factors are shard count and total checkpoint size. For \texttt{ServerlessLLM}, the key factors are total bytes, partition count, and the degree to which the workload is range-heavy. In broad terms, the design target is simple: smaller eager loads should stay on the simpler synchronous path, while larger loads should switch to \texttt{io\_uring}. The partitioned format adds nuance around this rule, which is why the selector remains format-aware rather than purely size-aware.

\section{Backend Designs}

\subsection{Synchronous Backend}

The synchronous backend uses blocking POSIX reads. In the final design it is not a purely serial implementation: large files are split into chunks and loaded with host-side thread parallelism. This keeps the backend simple at the API boundary while still exploiting multiple cores where appropriate. The design intent is not to make \texttt{sync} look trivial, but to give it a fair and practical implementation baseline.

\subsection{Async Backend}

The async backend is built on Tokio~\cite{tokio_docs,tokio_repo}. Its design is strongest for grouped and range-heavy workloads rather than for large whole-file eager loads. The final implementation groups work per file, uses dynamic batching, and reduces unnecessary task fanout. This backend remains important because the final results show that explicit async loading is often strong for \texttt{ServerlessLLM}-style partitioned workloads even when it is weak for \texttt{SafeTensors}.

\subsection{io\_uring Backend}

The final \texttt{io\_uring} backend is the most important architectural outcome of the project. The original prototype used one shared ring and one issuer thread. Profiling showed that this design badly underutilised the machine. The final backend instead uses multiple worker-owned rings, lazy ring initialisation, dynamic ring sizing, and workload-driven planning for chunk size and in-flight depth.

This is a crucial design point. The project does not conclude that \texttt{io\_uring} is inherently superior. It concludes that a well-designed \texttt{io\_uring} backend becomes attractive for larger workloads, whereas a poorly designed one does not.

\subsection{Memory-Mapped Path}

The mmap-backed path is retained as a useful lazy-access mode, especially for Python and vLLM experiments where a caller may not need to materialise every tensor immediately. It is not treated as the main eager-loading solution to the cold-start problem, because mmap solves a different problem: efficient selective access rather than aggressive eager loading.

\section{Checkpoint Formats}

TensorStore supports two checkpoint layouts.

\texttt{SafeTensors} stores tensors in a whole-file or multi-shard layout. Once the header is parsed, the loader sees a largely contiguous eager-loading workload. Figure~\ref{fig:safetensors_layout} shows the layout.

\begin{figure}[h]
    \centering
    \includegraphics[width=0.75\textwidth]{\figdir/safetensors-layout.png}
    \caption{SafeTensors file layout. The header describes each tensor and its byte offset within the file. Data is stored contiguously.}
    \label{fig:safetensors_layout}
\end{figure}

The \texttt{ServerlessLLM} layout distributes tensors across multiple partition files with a central index~\cite{FU:2024,serverlessllm_repo}. This introduces a more fragmented range-oriented workload and therefore stresses a different part of the backend design space. Figure~\ref{fig:serverlessllm_layout} shows the structure.

\begin{figure}[h]
    \centering
    \includegraphics[width=0.75\textwidth]{\figdir/serverlessllm-layout.png}
    \caption{ServerlessLLM partitioned checkpoint structure. Tensors are distributed across multiple partition files, with a central index mapping tensor names to locations.}
    \label{fig:serverlessllm_layout}
\end{figure}

The purpose of supporting both formats is not just compatibility. It is to expose two different I/O regimes under one framework.

\section{Python Bindings}

TensorStore exposes its functionality through Python bindings implemented with PyO3. The bindings intentionally expose a small public API: a \texttt{default} eager path, explicit \texttt{sync} and \texttt{async} eager paths, and lazy \texttt{open\_*} mmap-backed handles. On Linux, \texttt{io\_uring} remains an internal backend behind \texttt{default} in the public Python API, while benchmark-only explicit paths were used in the vLLM benchmark layer.

One important binding design decision is that the async Python path runs on a shared Tokio runtime rather than constructing a fresh runtime per call. This was one of the practical optimisations that materially improved Python-side performance.

\section{vLLM Integration}

The vLLM integration is implemented through a custom model-loader path. vLLM is particularly useful as an evaluation target because it separates model initialisation, time-to-first-token, and steady-state decode. This allows the report to test whether backend gains survive contact with a real serving engine.

The key design constraint is that vLLM's loader interface is synchronous. The final benchmark integration therefore compares synchronous and default eager-loading policies directly, and uses benchmark-only explicit loader variants where necessary to expose additional backend behaviour. This is a practical design choice: it reflects the serving interface that actually exists, rather than inventing an async serving path that the engine does not support.

\section{Cache-Control Design}

Cold-cache control is intentionally kept outside the core implementation. Earlier project iterations experimented with in-code cache eviction, but that approach was removed from the final system because it mixed benchmark ritual with loader logic. The final design instead treats cold-cache control as part of the experiment harness: the Linux shell is responsible for running \texttt{sync} and \texttt{drop\_caches} before benchmark invocations.

This decision is important for clarity. The loader implementation should express how TensorStore loads checkpoints; the benchmark harness should express how cold and warm conditions are created.

\section{Design Decisions and Trade-offs}

Table~\ref{tab:design_decisions} summarises the most important architectural decisions.

\begin{sidewaystable}
\small
\caption{Primary design decisions and trade-offs}
\label{tab:design_decisions}
\centering
\begin{tabular}{p{3cm}p{3.5cm}p{5cm}}
\toprule
\textbf{Decision} & \textbf{Alternative} & \textbf{Rationale} \\
\midrule
Rust core engine & Python or C++ & Rust provides memory safety without garbage collector overhead, direct access to low-level Linux \texttt{io\_uring}, and PyO3 bindings with minimal overhead. Python's GIL limits true parallelism; C++ lacks memory safety guarantees. \\
\midrule
Common loading interface & Backend-specific APIs & A shared interface allows fair backend comparison and lets the selector switch policies without rewriting format logic. \\
\midrule
Adaptive default policy & Fixed backend policy & Final evidence shows no universal winner; backend choice must depend on workload regime. \\
\midrule
Multi-worker \texttt{io\_uring} & Single shared ring & Profiling showed that a single-ring design underutilised the machine, while worker-owned rings materially improved performance. \\
\midrule
Shared Tokio runtime in bindings & Per-call runtime construction & A shared runtime reduces Python binding overhead and reflects how real async integrations are used. \\
\midrule
External cold-cache control & In-code cache eviction & Separating benchmark ritual from loader logic produces a cleaner implementation and clearer methodology. \\
\bottomrule
\end{tabular}
\normalsize
\end{sidewaystable}
